\section{Introduction}

The repository now contains enough ARC material that separate local papers are no longer the right endpoint. The LLM psychometric paper established that ARC can behave like a surprisingly clean measurement instrument for frontier language models. The human paper showed that sparse ARC-AGI-2 human logs still support a coherent latent difficulty structure and that human-vs-model alignment is neither trivial nor random. The non-LLM paper showed that current non-LLM systems are more mixed than hoped: they are not psychometrically empty, but they do not currently match the LLM consensus in reproducing the human difficulty profile. The solver-complexity package then changed the center of gravity of the project by showing that validated solver structure predicts LLM difficulty strongly and predicts human difficulty much more weakly. The cross-ARC package widened the matched human/LLM coverage and clarified what can and cannot be said about ARC-AGI-1 versus ARC-AGI-2. The efficiency package brought in performance/effort tradeoffs. And the external paper synthesis on URM, TRM, VARC, and McGovern suggested that recent non-LLM ARC progress may already share a deeper mechanistic loop.

Those results belong in one paper because they answer a common question from different angles: what kind of thing is ARC measuring across humans and machines, and what does the current evidence imply about the burdens each source family faces when solving it? If these workstreams are left isolated, several distortions become hard to avoid. The LLM psychometric story can start to sound stronger than it is about human comparability. The human paper can be read as if overlap tables alone settle the issue. The non-LLM paper can be read as either too pessimistic or too optimistic, depending on which nulls are emphasized. The solver-complexity results can look like an isolated curiosity rather than the integrative center of the repo. And the external solver papers can seem disconnected from the empirical evidence, even though they may actually be the best mechanistic explanation for several mixed results.

The purpose of this master paper is therefore twofold. First, it reconstructs the analytic pipeline at repo scale: what data exist, which units of analysis were actually used, where the overlaps are wide versus narrow, what inferential conventions were followed, and how each workstream depends on the others. Second, it offers a unified interpretation. The strongest version of that interpretation is that ARC contains a shared difficulty structure across humans and machine systems, but that shared structure is incomplete. LLM difficulty is especially sensitive to the structural complexity of the final successful solver. Human difficulty is less structure-loaded and more tied to search effort, time cost, and pair-level variation inside the same task. Current non-LLM systems are not yet broad human-like replicas, but they do supply complementary signal and point toward a mechanistic loop built from inductive bias, iterative refinement, and test-time candidate selection.

That framing also fits the benchmark's own original motivation. ARC was introduced as a measure of general fluid intelligence under severe sample constraints, not as another broad data-rich accuracy benchmark \cite{chollet2019arc}. ARC-AGI-2 tightened those constraints while shifting the difficulty distribution upward \cite{arcprizefoundation2025}. In parallel, recent non-LLM papers have argued, in different ways, that successful ARC solvers rely on stronger priors, more structured inference, and more deliberate test-time adaptation than standard large-model scaling stories would suggest \cite{gao2025urm, hu2025varc, jolicoeur2025trm, mcgovern2025ttatrm}. The repository now has enough internal evidence to evaluate those architectural claims against actual behavioral signatures.

The paper proceeds in nine large pieces. Section 2 lays out the repository-wide data and methods, with special emphasis on overlap structure and unit mismatch. Section 3 revisits the LLM psychometric line, including the broader benchmark-manifold addendum. Section 4 reconstructs the human psychometric case under sparse coverage. Section 5 addresses the non-LLM systems and explains why their result is simultaneously disappointing and still worth keeping. Section 6 treats solver complexity as the central integrative analysis layer. Section 7 adds the cross-ARC latent linkage and the efficiency results. Section 8 connects the empirical evidence to the external architecture papers. Section 9 states the strongest supported thesis, the main unresolved issues, and the practical next steps for turning this scaffold into a final long-form manuscript.
